\documentclass[en, color=black, device=pad, mode=geey]{elegantnote}

\usepackage{listings}
\usepackage{xcolor}
\usepackage[T1]{fontenc}
\usepackage{polski}

\title{Temat 2. Złożoność obliczeniowa }
\author{}
\date{\today}
\lstset{
    basicstyle=\ttfamily\small,
    keywordstyle=\color{blue},
    stringstyle=\color{green!50!black},
    commentstyle=\color{gray},
    numbers=left,
    numberstyle=\color{gray},
    stepnumber=1,
    frame=single,
    breaklines=true,
    captionpos=b
    showstringspaces=false, 
    showspaces=false
    commentstyle=\color{gray}\itshape,
    morecomment=[l]{//},
    literate={ą}{{\k{a}}}1 {ć}{{\'c}}1 {ę}{{\k{e}}}1 {ł}{{\l{}}}1 {ń}{{\'n}}1 {ó}{{\'o}}1 {ś}{{\'s}}1 {ź}{{\'z}}1 {ż}{{\.z}}1
             {Ą}{{\k{A}}}1 {Ć}{{\'C}}1 {Ę}{{\k{E}}}1 {Ł}{{\L{}}}1 {Ń}{{\'N}}1 {Ó}{{\'O}}1 {Ś}{{\'S}}1 {Ź}{{\'Z}}1 {Ż}{{\.Z}}1,
    keywordstyle=\color{blue}\bfseries,
    morekeywords={READ, LOAD, STORE, JZERO, SUB, JGTZ, JUMP, WRITE, HALT}, 
}

\begin{document}

\maketitle


\newpage

\section{Zadanie 1}

\subsection{}
Rejestry: $R_1$ będzie przechowywać bieżącą wartość $x$, $R_2$ będzie przechowywać dotychczasowe $max$. Akumulator to $R_0$.


\begin{lstlisting}
READ 1      // czytaj x do R1
LOAD 1      // wczytaj x do akumulatora
STORE 2     // max <- x (zapisz w R2)
LOOP: LOAD 1 // wczytaj x do sprawdzania warunku pętli
JZERO END   // jeśli x == 0, zakończ pętlę i skocz do END
SUB 2       // Akumulator = x - max
JGTZ UPDATE // jeśli x - max > 0 (czyli x max), skocz do UPDATE
JUMP NEXT   // w przeciwnym razie omiń aktualizację
UPDATE: LOAD 1 // wczytaj x
STORE 2     // max <- x
NEXT: READ 1 // czytaj kolejne x
JUMP LOOP   // wróć na początek pętli
END: WRITE 2 // wypisz max (z R2)
HALT        // koniec programu
\end{lstlisting}

\newpage
\subsection{Złożoność czasowa i pamięciowa: }


Niech:
\begin{itemize}
    \item  $N$ - oznacza liczbę wszystkich elementów ciągu wejściowego (łącznie z kończącym go zerem).
    \item $V$ - oznacza maksymalną wartość bezwzględną liczby w ciągu ($V = \max(|x_i|)$).
    \item $l(x)$ - oznacza długość bitową (koszt logarytmiczny) liczby $x$, standardowo $l(x) = O(\log |x|)$.
\end{itemize}

\subsubsection{Kryterium równomierne (jednorodne)}
W tym modelu zakładamy, że każda instrukcja wykonuje się w czasie $O(1)$, a każdy rejestr zajmuje pamięć $O(1)$, niezależnie od wielkości przechowywanych w nim liczb.

\begin{itemize}
    \item Złożoność pamięciowa, $S(N)$: Program potrzebuje zaledwie kilku rejestrów do działania niezależnie od rozmiaru danych wejściowych: akumulatora na operacje (zwykle $R_0$), rejestru na bieżącą wartość $x$ (np. $R_1$) oraz rejestru na dotychczasowe maksimum $max$ (np. $R_2$). Liczba zajmowanych komórek pamięci jest stała.
          \begin{equation}
              S(N) = O(1)
          \end{equation}
    \item Złożoność czasowa, $T(N)$: Pętla dopóki wykona się dokładnie $N-1$ razy (dla każdej liczby aż do napotkania zera). Wewnątrz pętli znajduje się stała liczba prostych instrukcji (porównanie, ewentualne przypisanie, wczytanie kolejnej liczby). Wykonanie każdej z nich kosztuje $O(1)$.
          \begin{equation}
              T(N) = O(N)
          \end{equation}
\end{itemize}


\subsubsection{Kryterium logarytmiczne}
W tym modelu koszt wykonania instrukcji oraz koszt przechowania wartości zależy od długości binarnej operanda. Koszt przechowania i operowania na liczbie $x$ wynosi $l(x) \approx O(\log |x|)$.

\begin{itemize}
    \item Złożoność pamięciowa, S(N): Wykorzystujemy stałą liczbę rejestrów, ale ich rozmiar zależy od tego, jak duże liczby w nich przechowujemy. Największą liczbą, jaka może znaleźć się w pamięci (w zmiennej x lub max), jest V. Koszt zapamiętania takiej liczby to O(logV).
          \begin{equation}
              S(N) = O(\log V)
          \end{equation}
    \item Złożoność czasowa, $T(N)$: Pętla wykonuje się $N-1$ razy. W każdym obrocie pętli wykonujemy operacje (np. wczytywanie READ, odejmowanie przy warunku jeśli x>max, zapisywanie STORE). W najgorszym przypadku w każdym kroku operujemy na liczbach rzędu wielkości $V$. Koszt pojedynczej operacji na takich liczbach wynosi $O(\log V)$. Sumując ten koszt dla $N$ obrotów pętli, otrzymujemy:
          \begin{equation}
              T(N) = O(N \log V)
          \end{equation}
\end{itemize}

\newpage

\section{Zadanie 2}
Algorytm zlicza wystąpienia jedynek i dwójek, a na końcu - po napotkaniu zera - porównuje oba liczniki.

\begin{lstlisting}
    c1 <- 0
    c2 <- 0
    czytaj x
    dopóki x != 0 wykonuj {
        jeśli x == 1 to c1 <- c1 + 1
        jeśli x == 2 to c2 <- c2 + 1
        czytaj x
    }
    jeśli c1 == c2 to pisz 1
\end{lstlisting}

\newpage
\subsection{Program w języku maszyny RAM}
Przyjmuję następujące przeznaczenie rejestrów: $R_1$ - przechowuje aktualnie wczytaną wartość $x$
$R_2$ - licznik jedynek ($c_1$)
$R_3$ - licznik dwójek ($c_2$)
$R_0$ - akumulator

\begin{lstlisting}
LOAD =0      // przygotuj 0 w akumulatorze
STORE 2      // c1 = 0
STORE 3      // c2 = 0
LOOP: READ 1  // czytaj kolejne x do R1
LOAD 1       // wczytaj x do akumulatora
JZERO CHECK  // jeśli x == 0, zakończ pętlę i sprawdź liczniki
SUB =1       // Akumulator = x - 1
JZERO IS_ONE // jeśli wynik to 0 (czyli x == 1), skok
LOAD 1       // wczytaj ponownie x do akumulatora
SUB =2       // Akumulator = x - 2
JZERO IS_TWO // jeśli wynik to 0 (czyli x == 2), skok
JUMP LOOP // (wraca na początek jeśli to coś innego)
IS_ONE: LOAD 2 // wczytaj licznik c1
ADD =1       // c1 + 1
STORE 2      // zapisz zaktualizowane c1
JUMP LOOP    // wróć po kolejną liczbę
IS_TWO: LOAD 3 // wczytaj licznik c2
ADD =1       // c2 + 1
STORE 3      // zapisz zaktualizowane c2
JUMP LOOP    // wróć po kolejną liczbę
CHECK: LOAD 2 // wczytaj c1
SUB 3        // Akumulator = c1 - c2
JZERO EQUAL  // jeśli różnica wynosi 0, są równe
HALT         // jeśli nierówne, po prostu zatrzymaj program
EQUAL: WRITE =1 // wypisz 1 (skoro są równe)
HALT         // zatrzymaj program
\end{lstlisting}

\subsection{Złożoność czasowa}

3. Złożoność czasowa i pamięciowa (podpunkt b)Niech $N$ oznacza całkowitą liczbę elementów podanych na wejściu (wliczając kończące zero). Kryterium równomierne (jednorodne)Zakładamy, że każda instrukcja oraz każda komórka pamięci kosztuje $O(1)$.Złożoność pamięciowa: Wykorzystujemy stałą liczbę rejestrów ($R_0, R_1, R_2, R_3$), niezależnie od długości ciągu na wejściu.$$S(N) = O(1)$$Złożoność czasowa: Pętla odczytująca wykonuje się $N$ razy. W każdej iteracji maszyna wykonuje co najwyżej kilka instrukcji (wczytanie, 1-2 porównania, dodanie do licznika, skok). Każda z nich kosztuje czas stały.$$T(N) = O(N)$$Kryterium logarytmiczneZakładamy, że koszt zależy od długości bitowej przetwarzanych i przechowywanych wartości ($l(x) \approx O(\log |x|)$).Złożoność pamięciowa: Same liczby z wejścia (0, 1, 2) zajmują mało miejsca, ale liczniki $c_1$ i $c_2$ mogą urosnąć nawet do wartości $N$ (jeśli ciąg składa się z samych jedynek albo samych dwójek). Aby przechować w rejestrze wartość $N$, potrzebujemy logarytmicznej liczby bitów.$$S(N) = O(\log N)$$Złożoność czasowa: W każdym kroku pętli wykonujemy dodawanie do licznika. Dodanie wartości 1 do licznika przechowującego w danym momencie liczbę $k$ kosztuje czas rzędu $O(\log k)$ (ponieważ musimy przetworzyć wszystkie bity tej liczby). Sumując te koszty dla wszystkich $N$ elementów, otrzymujemy ciąg sumaryczny: $\sum_{k=1}^{N} O(\log k)$. Ze wzoru Stirlinga (lub analizy matematycznej na logarytmach) wiemy, że $\log(N!) \approx N \log N$.$$T(N) = O(N \log N)$$


\newpage

\section{Napisać w pseudojęzyku algorytm wyznaczania wartości n!
  dla podanej liczby nieujemne}

Zaczynamy od napisania prostego algorytmu iteracyjnego. Przypisujemy wynikowi początkową wartość 1, a następnie w pętli mnożymy go przez kolejne wartości $n$, po czym zmniejszamy $n$ o 1, aż dojdziemy do zera.

\begin{lstlisting}
czytaj n
silnia <- 1
dopóki n > 0 wykonuj {
    silnia <- silnia * n
    n <- n - 1
}
pisz silnia
\end{lstlisting}

\newpage
\subsection{Program w języku maszyny RAM}
Przydział rejestrów:$R_1$ – przechowuje zmienną $n$$R_2$ – przechowuje aktualnie obliczoną wartość $silnia$$R_0$ – akumulator (używany domyślnie do operacji)

        \begin{lstlisting}
READ 1      // wczytaj wartość n do R1
LOAD =1     // załaduj stałą 1 do akumulatora
STORE 2     // silnia = 1 (zapisz początkową wartość do R2)
LOOP: LOAD 1 // wczytaj aktualne n z R1 do akumulatora
JZERO END   // jeśli n == 0, zakończ pętlę i przejdź do END
MULT 2      // Akumulator = n * silnia
STORE 2     // silnia = n * silnia (zapisz nowy wynik w R2)
LOAD 1      // wczytaj n z R1 do akumulatora
SUB =1      // Akumulator = n - 1
STORE 1     // n = n - 1 (zapisz zmniejszone n w R1)
JUMP LOOP   // wróć na początek pętli
END: WRITE 2 // wypisz wynik (wartość z R2)
HALT        // zatrzymaj program
\end{lstlisting}

        \newpage
        \subsection{Złożoność czasowa i pamięciowa}
        Kryterium równomierne (jednorodne)W tym modelu każda instrukcja oraz każda użyta komórka pamięci ma stały koszt, niezależnie od wielkości przetwarzanych liczb.Złożoność pamięciowa: Program korzysta tylko z trzech rejestrów ($R_0$, $R_1$, $R_2$). Ich liczba nie zmienia się wraz ze wzrostem $n$.$$S(n) = O(1)$$Złożoność czasowa: Pętla wykonuje się dokładnie $n$ razy. W każdym kroku pętli wykonywana jest stała liczba prostych instrukcji.$$T(n) = O(n)$$Kryterium logarytmiczneW tym modelu koszt przechowania liczby oraz operacji na niej jest proporcjonalny do jej długości bitowej, czyli zależy od wartości logarytmu z tej liczby ($l(x) = O(\log |x|)$).Złożoność pamięciowa: W rejestrze $R_2$ na samym końcu przechowywana jest wartość $n!$. Ze wzoru Stirlinga wiemy, że długość zapisu binarnego liczby $n!$ wynosi w przybliżeniu $n \log n$. Zatem liczba bitów (pamięć) potrzebna do jej zapisania rośnie nieliniowo.$$S(n) = O(n \log n)$$Złożoność czasowa: W $i$-tym kroku pętli maszyna wykonuje mnożenie (instrukcja MULT) liczby $i$ (o długości $O(\log i)$) przez dotychczasowy iloczyn $i!$ (o długości $O(i \log i)$). Koszt takiej instrukcji w modelu logarytmicznym to suma długości bitowych argumentów, czyli $O(i \log i)$. Aby obliczyć całkowity czas, musimy zsumować ten koszt dla wszystkich $n$ kroków pętli:$$T(n) = \sum_{i=1}^{n} O(i \log i)$$Z własności sumowania ciągów otrzymujemy ostateczną złożoność:$$T(n) = O(n^2 \log n)$$
\end{document}